
\documentclass{article}
\usepackage{colortbl}
\usepackage{makecell}
\usepackage{multirow}
\usepackage{supertabular}

\begin{document}

\newcounter{utterance}

\centering \large Interaction Transcript for game `cladder', experiment `full\_v1.5\_default', episode 687 with Qwen3.5{-}9B.
\vspace{24pt}

{ \footnotesize  \setcounter{utterance}{1}
\setlength{\tabcolsep}{0pt}
\begin{supertabular}{c@{$\;$}|p{.15\linewidth}@{}p{.15\linewidth}p{.15\linewidth}p{.15\linewidth}p{.15\linewidth}p{.15\linewidth}}
    \# & \multicolumn{2}{c}{Player} && \multicolumn{2}{c}{Game Master} \\
    \hline

    \theutterance \stepcounter{utterance}  
    & & & \multicolumn{4}{p{0.6\linewidth}}{
        \cellcolor[rgb]{0.9,0.9,0.9}{
            \makecell[{{p{\linewidth}}}]{
                \texttt{\tiny{[P1$\langle$GM]}}
                \texttt{You are an expert in causal inference. The following question is not a typical commonsense query, but rather a meticulously designed question created by a professor specializing in causal inference, intended to assess the students' mastery of the course content.} \\
\texttt{Imagine a self{-}contained, hypothetical world with only the following conditions, and without any unmentioned factors or causal relationships: Gender has a direct effect on department competitiveness and admission status. Department competitiveness has a direct effect on admission status. For individuals who are not male and applicants to a non{-}competitive department, the probability of admission acceptance is 58\%. For individuals who are not male and applicants to a competitive department, the probability of admission acceptance is 25\%. For individuals who are male and applicants to a non{-}competitive department, the probability of admission acceptance is 59\%. For individuals who are male and applicants to a competitive department, the probability of admission acceptance is 28\%. For individuals who are not male, the probability of competitive department is 29\%. For individuals who are male, the probability of competitive department is 77\%. Does gender positively affect admission status through department competitiveness?} \\
\texttt{Start your answer with 'Yes' or 'No', followed by additional reasoning or evidence to support your explanation.} \\
            }
        }
    }
     \\ \\

    \theutterance \stepcounter{utterance}  
    & \multicolumn{4}{p{0.6\linewidth}}{
        \cellcolor[rgb]{0.561,0.737,0.561}{
            \makecell[{{p{\linewidth}}}]{
                \texttt{\tiny{[P1$\rangle$GM]}}
                \texttt{No} \\
\\ 
\texttt{**Reasoning:**} \\
\\ 
\texttt{To determine if gender positively affects admission status *through* department competitiveness, we must calculate the **Natural Direct Effect (NDE)** and the **Natural Indirect Effect (NIE)** of changing gender from "not male" to "male," while holding the mediator (department competitiveness) constant at its value under the "not male" condition.} \\
\\ 
\texttt{Let:} \\
\texttt{*   \$G\$ = Gender (0 = not male, 1 = male)} \\
\texttt{*   \$D\$ = Department Competitiveness (0 = non{-}competitive, 1 = competitive)} \\
\texttt{*   \$Y\$ = Admission Status (0 = rejected, 1 = accepted)} \\
\\ 
\texttt{**1. Identify the Probabilities:**} \\
\texttt{*   \$P(D=1 | G=0) = 0.29\$ (Probability of competitive dept for not male)} \\
\texttt{*   \$P(D=1 | G=1) = 0.77\$ (Probability of competitive dept for male)} \\
\texttt{*   \$P(Y=1 | G=0, D=0) = 0.58\$} \\
\texttt{*   \$P(Y=1 | G=0, D=1) = 0.25\$} \\
\texttt{*   \$P(Y=1 | G=1, D=0) = 0.59\$} \\
\texttt{*   \$P(Y=1 | G=1, D=1) = 0.2} \\
            }
        }
    }
    & & \\ \\

    \theutterance \stepcounter{utterance}  
    & & & \multicolumn{2}{p{0.3\linewidth}}{
        \cellcolor[rgb]{0.9,0.9,0.9}{
            \makecell[{{p{\linewidth}}}]{
                \texttt{\tiny{[GM$|$GM]}}
                \texttt{No} \\
            }
        }
    }
    & & \\ \\

    \theutterance \stepcounter{utterance}  
    & & & \multicolumn{2}{p{0.3\linewidth}}{
        \cellcolor[rgb]{0.9,0.9,0.9}{
            \makecell[{{p{\linewidth}}}]{
                \texttt{\tiny{[GM$|$GM]}}
                \texttt{no} \\
            }
        }
    }
    & & \\ \\

    \theutterance \stepcounter{utterance}  
    & & & \multicolumn{2}{p{0.3\linewidth}}{
        \cellcolor[rgb]{0.9,0.9,0.9}{
            \makecell[{{p{\linewidth}}}]{
                \texttt{\tiny{[GM$|$GM]}}
                \texttt{game\_result = WIN} \\
            }
        }
    }
    & & \\ \\

\end{supertabular}
}

\end{document}
